\section{Model Architecture}

The primary architecture used in this study is ResNet-18, based on the residual learning
framework of He et al.~\cite{he2016deep}, adapted for CIFAR-10 resolution. ResNet-18
provides a useful balance between expressiveness and computational cost: it is large
enough to support nontrivial adversarial training experiments, but still efficient enough
to run repeated-seed comparisons and ablation studies.

The classifier is denoted by \(f_\theta\). For an input image \(x\), the model outputs logits
$$
z = f_\theta(x) \in \mathbb{R}^{C},
$$
where \(C=10\) is the number of CIFAR-10 classes. In addition to logits, the proposed method requires access to the penultimate-layer representation
$$
h = \phi_\theta(x) \in \mathbb{R}^{d_h}.
$$

This feature vector is used for clustering adversarial examples in AttackDRO++ and for constructing the gradient-fingerprint augmented feature described in Chapter 4. For ResNet-18 on CIFAR-10, the penultimate representation dimension is \(d_h = 512\). Therefore, the raw classifier-head gradient proxy has dimension \(C d_h = 10 \times 512 = 5120\), which is later reduced by random projection.

Unless otherwise stated, all main experiments use the same ResNet-18 architecture and differ only in the training objective. This controlled setup ensures that performance differences can be attributed to the training method rather than to architectural changes. Higher-capacity architectures are treated as scaling checks or future extensions rather than the main basis for comparison.
